\scanpage{228}
\mspage{178}
% L1–2
induit un iso pour la cohomologie à \unsure{cft.|cfft.} loc. constants, jusqu'en dimension $n-1$ inclus.

% L3–6
\deleted{\diagram{\begin{tikzcd}[ampersand replacement=\&,column sep=large]
E(X)=\unsure{E|F}(X,x)=\mathrm{Ch}(X,x,-)\arrow[r]\arrow[d] \& E(Y,y)=E(Y)\arrow[d]\\
X\arrow[r,"f"] \& Y\\
f^*F=F_X \& F
\end{tikzcd}}}
% L7 left
\deleted{\[
H^*(E(X)\unsure{,|)}
\]}

% L7–8
Finalement, je me suis $\pm$ convaincu que l'énoncé du th. 1 est faux tel quelle.
% L9–14
Voici un énoncé que j'espère vrai, et que je vais essayer de prouver dans un esprit ``utilitaire'', en utilisant tout ce que je sais en homotopie et cohomologie (et pas seulement ce que j'ai développé ``from scratch'' dans le modélisateur $\mathrm{Cat}$).

% L15–20
\uline{Théorème 1 rectifié} Considérons pour un entier $n\deleted{\unsure{>}}\geq2$ donné, les conditions suivantes sur une flèche $f:X\to Y$ de $\mathrm{Cat}$ (ou une application continue d'espaces \unsure{topologiques|top. logiques}, \deleted{\unsure{x}} ou un morphisme dans $\widehat{\Delta}$ etc etc) :
% L21–23
\deleted{\unsure{a}} $a_n$) Si $Z$ est une ``fibre homotopique'' de $f$ (NB il y en a essentiellement une par composante connexe de $Y$), alors \deleted{\unsure{$f$}} $Z$ est $n$-connexe, i.e.
% archive sheet number: 228
